\documentclass{article}

\usepackage{amsmath}
\usepackage{amssymb}
\usepackage{geometry}
\geometry{margin=2.5cm}
\title{Tensor Omics: Data Model Definitions}
\author{Asis Hallab, ``el Bobito''}

\begin{document}

\maketitle

\section{Data Models}


All data in Tensor Omics (TOX) must be stored continuously in memory in basic
Fortran types. All data must be serializable and deserializable from binary
files (\texttt{.txdata} format). Those files consist of a header section
informing about the contained basic Fortran structures (scalars and arrays),
their sizes, and the offset of where to start to read them.

Tensor Omics will require meta-data to be readable and writable from data
tables, e.g. \texttt{CSV} format. To represent this efficiently we will use the
below \texttt{DataTable} definitions.

\subsection{DataTable}

\subsubsection{Overview}

The \textbf{F42 DataTable} implements the \textbf{F42 DataStructure} standard.
It provides a compact, efficient, and portable representation of structured tabular data (e.g., parsed from CSV files).
It is optimized for serialization, patching, and memory-constrained environments.

\subsubsection{Public Methods}

\begin{itemize}
    \item \texttt{read\_from\_file}(file, sep\_char, column\_types, max\_char\_length, n\_rows, column\_names)
    \item \texttt{calculate\_memory\_requirements}()
    \item \texttt{serialize}()
    \item \texttt{deserialize}()
    \item \texttt{log}()
    \item \texttt{update()} \hfill (Detailed definition later)
    \item \texttt{save\_patch()} \hfill (Detailed definition later)
\end{itemize}

\subsubsection{Internal Structure}

The DataTable is composed of the following internal components:

\begin{itemize}
    \item \textbf{Type--Mapping Array}: a \( 2 \times k \) integer array.
    \begin{itemize}
        \item Row 1: Stores the column type for each input column.
        \item Row 2: Stores the index of the corresponding column inside the appropriate data container (see below).
    \end{itemize}
    \item \textbf{Data Containers}:
    \begin{itemize}
        \item \texttt{int\_array}: A \( n\_rows \times i \) array for integer columns.
        \item \texttt{real\_array}: A \( n\_rows \times r \) array for real (floating-point) columns.
        \item \texttt{char\_array}: A \( n\_rows \times c \) array of fixed-length character sequences (length = max\_char\_length) for string columns.
    \end{itemize}
    where \( i + r + c = k \), and each value \( i, r, c \geq 0 \).
    
    \item \textbf{Column Names}:
    \begin{itemize}
        \item A \( k \times 256 \) character array.
        \item Each entry holds the name of a column, truncated or padded to a maximum length of 256 characters.
    \end{itemize}
\end{itemize}

\subsubsection{Detailed Method Definitions}

\subsubsection*{read\_from\_file}%

Method signature: \texttt{read\_from\_file(file, sep\_char, column\_types, max\_char\_length, n\_rows, column\_names)}

Reads tabular data from a file and populates the internal structure.

\begin{itemize}
    \item \textbf{file}: Path or handle of the input file.
    \item \textbf{sep\_char}: Character used to separate fields (e.g., `,` or `;`).
    \item \textbf{column\_types}: Integer array of length \( k \) specifying the type of each input column:
    \begin{itemize}
        \item 1 = Integer
        \item 2 = Real (floating point)
        \item 3 = Character (string)
    \end{itemize}
    \item \textbf{max\_char\_length}: Maximum allowed character length for each string field (in char columns).
    \item \textbf{n\_rows}: Expected number of data rows (excluding headers).
    \item \textbf{column\_names} (optional): Array of column names. If omitted, defaults to generic names ("col1", "col2", ...).
\end{itemize}

\subsubsection*{calculate\_memory\_requirements}

Method signature: \texttt{calculate\_memory\_requirements(n\_rows, col\_types)}

Estimates and returns the total memory (in bytes) required for the DataTable, including:

\begin{itemize}
    \item Size of Type--Mapping Array: \( 2 \times k \times \text{sizeof(int)} \)
    \item Size of Data Containers:
    \begin{itemize}
        \item \( n\_rows \times i \times \text{sizeof(int)} \)
        \item \( n\_rows \times r \times \text{sizeof(real)} \)
        \item \( n\_rows \times c \times \text{max\_char\_length} \times \text{sizeof(char)} \)
    \end{itemize}
    \item Size of Column Names Array: \( k \times 256 \times \text{sizeof(char)} \)
\end{itemize}

\subsubsection*{serialize()}

Serializes the DataTable into a binary format consisting of:

\begin{enumerate}
    \item Header block (generated by \texttt{generate\_header}).
    \item Raw memory blocks for:
    \begin{itemize}
        \item Type--Mapping Array
        \item int\_array
        \item real\_array
        \item char\_array
        \item column\_names
    \end{itemize}
\end{enumerate}

All blocks are concatenated sequentially without gaps.
Serialization follows little-endian ordering unless otherwise specified.

\subsubsection*{generate\_header()}

Creates a header structure containing metadata such as:

\begin{itemize}
    \item Magic number (identifies the file type, e.g., 'F42DTBL')
    \item Version number
    \item Number of rows (\( n\_rows \))
    \item Number of columns (\( k \))
    \item Maximum character length (\( max\_char\_length \))
    \item Column counts per type (\( i, r, c \))
\end{itemize}

\subsubsection*{deserialize()}

Inverse operation of \texttt{serialize()}:

\begin{itemize}
    \item Reads and verifies the header.
    \item Reconstructs the Type--Mapping Array and Data Containers from the binary block.
    \item Ensures compatibility of all dimensions and memory alignment.
\end{itemize}

\subsubsection{Methods to be Specified Later}

The following methods will be specified in detail in a future document:

\begin{itemize}
    \item \texttt{log()}
    \item \texttt{update()}
    \item \texttt{save\_patch()}
\end{itemize}

\subsubsection{Summary Diagram}

\[
\text{Input CSV}
\quad \xrightarrow{\text{read\_from\_file}} \quad
\left\{
\begin{array}{l}
\text{Type--Mapping Array} \\
\text{int\_array} \\
\text{real\_array} \\
\text{char\_array} \\
\text{column\_names}
\end{array}
\right\}
\quad \xrightarrow{\text{serialize}} \quad
\text{Binary Format}
\]

\section{F42-Compliant Nested HashMap System}

\subsection{Overview}

The F42 HashMap system provides an efficient, statically typed representation of arbitrarily nested key-value maps using flat, preallocated memory bands. Each scalar type is stored in a dedicated array (Turing Band). Keys and properties are indexed via integer references, enabling fast lookup, nesting, and serialization.

\subsection{Memory Bands and Data Layout}

\begin{itemize}
  \item \textbf{Value Bands} (type-specific storage):
    \begin{itemize}
      \item \texttt{int\_vals(n\_int)}
      \item \texttt{real\_vals(n\_real)}
      \item \texttt{char\_vals(n\_char, len\_char)}
    \end{itemize}

  \item \textbf{Key Index Table}:
    \[
    \texttt{keys(3, n\_maps)} \quad \text{contains for each HashMap entry:}
    \]
    \begin{itemize}
      \item \texttt{keys(1,i)}: Index of the key string in \texttt{char\_vals}
      \item \texttt{keys(2,i)}: Start index in \texttt{properties} table
      \item \texttt{keys(3,i)}: Stop index in \texttt{properties} table
    \end{itemize}

  \item \textbf{Properties Table}:
    \[
    \texttt{properties(4, n\_properties)}
    \]
    Each column corresponds to a field in a map:
    \begin{itemize}
      \item \texttt{properties(1,j)}: Field name (index in \texttt{char\_vals})
      \item \texttt{properties(2,j)}: Type code (1 = int, 2 = char, 3 = real, 4 = nested map)
      \item \texttt{properties(3,j)}: Start index in the corresponding value band
      \item \texttt{properties(4,j)}: Stop index (for vector values or nested maps)
    \end{itemize}
\end{itemize}

\subsection{Access Function Pseudocode}

\paragraph{Lookup Key in HashMap}
\begin{verbatim}
function get_map_range(keys, key_id):
  do i = 1, size(keys, 2)
    if keys(1, i) == key_id then
      return keys(2, i), keys(3, i) ! Return property start/stop indices
    end if
  end do
  return -1, -1 ! Not found
\end{verbatim}

\paragraph{Lookup Property in Map}
\begin{verbatim}
function get_property(properties, i_start, i_stop, field_id):
  do j = i_start, i_stop
    if properties(1, j) == field_id then
      type_id = properties(2, j)
      val_start = properties(3, j)
      val_stop  = properties(4, j)
      return type_id, val_start, val_stop
    end if
  end do
  return -1, -1, -1 ! Not found
\end{verbatim}

\paragraph{Retrieve Value}
Once type and start/stop indices are known:
\begin{verbatim}
if type_id == 1:
  return int_vals(val_start:val_stop)
if type_id == 2:
  return char_vals(val_start:val_stop, :)
...
if type_id == 4:
  return nested map keys(:, val_start:val_stop)
\end{verbatim}

\subsection{Serialization and Nested Maps}

\begin{itemize}
  \item Entire structure can be written as binary blocks:
    \begin{itemize}
      \item \texttt{keys}, \texttt{properties}, value bands
    \end{itemize}
  \item Nested maps are supported via:
    \begin{itemize}
      \item \texttt{type\_id = 4}
      \item \texttt{val\_start}, \texttt{val\_stop} refer to rows in \texttt{keys()}
    \end{itemize}
  \item Recursive traversal follows key -> property -> nested map pattern.
\end{itemize}

\subsection{Advantages}

\begin{itemize}
  \item Fully F42-compliant: flat arrays, preallocated memory, no dynamic allocation.
  \item Serialization-ready.
  \item Recursive, nested structure without pointers.
  \item Clean and performant traversal and update patterns.
\end{itemize}

\subsection{Test Cases for F42 Nested HashMaps}

\paragraph{Test 1: Simple Flat Key-Value Retrieval}

\begin{verbatim}
! Setup:
! char_vals: ["name", "age"]
! int_vals: [42]
! keys(1,1) = index("patient1") in char_vals
! keys(2,1:3,1) = (1,2)        ! property range

! properties(:,1) = [index("name"), 2, 1, 1]
! properties(:,2) = [index("age"),  1, 1, 1]

! Test retrieval of age
call get_map_range(keys, index("patient1")) => (1, 2)
call get_property(properties, 1, 2, index("age")) => (1, 1, 1)
int_vals(1) = 42 ! -> OK
\end{verbatim}

\paragraph{Test 2: Nested Map Traversal}

\begin{verbatim}
! Setup:
! keys(1,1) = index("patient1")
! properties(:,2) = [index("clinical_data"), 4, 2, 2]
! keys(:,2) = nested map root: index("diagnosis"), range (3,3)
! properties(:,3) = [index("diagnosis"), 2, 3, 3]
! char_vals(3,:) = "healthy"

! Access:
call get_map_range(keys, index("patient1")) => (1, 2)
call get_property(properties, 1, 2, index("clinical_data")) => (4, 2, 2)
call get_map_range(keys, 2) => (3, 3)
call get_property(properties, 3, 3, index("diagnosis")) => (2, 3, 3)
char_vals(3,:) = "healthy" ! -> OK
\end{verbatim}

\paragraph{Test 3: Missing Key or Property Handling}

\begin{verbatim}
! Try to get nonexistent key "ghost"
call get_map_range(keys, index("ghost")) => (-1, -1) => OK

! Try to get nonexistent field "weight" in patient1
call get_map_range(keys, index("patient1")) => (1,2)
call get_property(properties, 1, 2, index("weight")) => (-1, -1, -1) => OK
\end{verbatim}

\paragraph{Test 4: Multiple Fields of Same Type}

\begin{verbatim}
! Add "height" and "weight" of type real
! properties(:,4) = [index("height"), 3, 1, 1]
! properties(:,5) = [index("weight"), 3, 2, 2]
! real_vals(1) = 1.75
! real_vals(2) = 70.0

! Test:
call get_property(properties, 1, 5, index("weight")) => (3, 2, 2)
real_vals(2) = 70.0 => OK
\end{verbatim}

\paragraph{Test 5: Field as Vector Value}

\begin{verbatim}
! properties(:,6) = [index("expression"), 3, 3, 5]
! real_vals(3:5) = [1.2, 0.9, 3.1]

! Retrieve vector value
call get_property(properties, 1, 6, index("expression")) => (3, 3, 5)
real_vals(3:5) = [1.2, 0.9, 3.1] => OK
\end{verbatim}

\paragraph{Validation Summary}
Each of these tests ensures:
\begin{itemize}
  \item Correct indexing of key and property arrays
  \item Accurate lookup and retrieval in Turing Bands
  \item Support for scalar and vector fields
  \item Support for nested maps via integer nesting
  \item Clean handling of missing keys or fields
\end{itemize}

\noindent
Further tests may involve performance profiling, binary serialization
validation, and nested structure export/reconstruction checks.

\section{F42 Hashmap Memory and Update Semantics (Extended)}

\subsection{Compaction and Memory Hygiene}

To maintain efficient memory use and avoid fragmentation due to property overwrites or deletions, we define an explicit compaction mechanism for all Turing bands and relevant 2D structures such as the `keys` and `properties` tables.

\subsubsection{Compaction Strategy}

\begin{itemize}
  \item For each band (e.g., \texttt{int\_vals}, \texttt{char\_vals}, \texttt{real\_vals}), maintain an accompanying integer array \texttt{deleted\_indices(:)} containing the indices of entries marked as obsolete.
  \item During compaction:
    \begin{enumerate}
      \item Allocate a new compacted array of minimal size.
      \item Copy only non-obsolete entries.
      \item Record an index remapping from old to new positions.
      \item Traverse the \texttt{properties} array and remap all band references accordingly.
      \item Replace the old band with the new one.
    \end{enumerate}
  \item This strategy supports all scalar types and nested \texttt{HashMap} structures uniformly.
  \item \textbf{Do not embed deletion markers in the bands themselves}, as that strategy fails for fixed-length character arrays and would violate F42 separation-of-concerns principles.
  \item Logical deletion arrays also permit filtered read-only access (e.g., using \texttt{WHERE (.NOT. is\_deleted)}) without triggering compaction.
\end{itemize}

\subsubsection{HashMap Lifecycle Pseudocode (extended)}

\begin{verbatim}
subroutine compact_band(old_band, deleted_indices, new_band, index_map)
  integer, intent(in) :: old_band(:), deleted_indices(:)
  integer, intent(out) :: new_band(:), index_map(size(old_band))

  ! Step 1: mark deletions
  logical :: keep(size(old_band))
  keep = .true.
  keep(deleted_indices) = .false.

  ! Step 2: create new_band and remap indices
  new_band = pack(old_band, keep)
  new_pos = 1
  do i = 1, size(old_band)
    if (keep(i)) then
      index_map(i) = new_pos
      new_pos = new_pos + 1
    else
      index_map(i) = -1
    end if
  end do

  ! Step 3: update property lookup intervals using index_map
end subroutine

type(HashMapCollection) :: mapset
call compact_band(mapset%int_vals, mapset%int_deleted, compacted, remap)

doit: do i = 1, mapset%num_properties
  if (mapset%properties(2, i) == 1) then
    mapset%properties(3:4, i) = remap(mapset%properties(3:4, i))
  end if
end do doit
\end{verbatim}

This strategy provides robust, scalable memory hygiene compatible with both
Fortran memory models and F42 minimalistic reproducibility.

\subsection{F42-Compliant Graph Representation Using Flat Arrays}

\subsubsection{Design Overview}

The F42 graph structure is implemented using a pair of flat 2D arrays to represent connectivity and node/edge labeling. This design preserves strict memory control, serialization efficiency, and WebAssembly compatibility.

\begin{itemize}
  \item \textbf{Adjacency Matrix} \quad \texttt{adj(n, m)} \\
    Each entry at \texttt{adj(i, j)} holds the label index (column in the label matrix) of the node to which node \texttt{i} is connected via edge \texttt{j}.
  \item \textbf{Label Matrix} \quad \texttt{labels(2, k)} \\
    - Row 1: Label type.\\
    \hspace*{1.5em} \texttt{1 = DataTable row} \\
    \hspace*{1.5em} \texttt{2 = HashMap reference} \\
    - Row 2: Label index (row or key index).

    The first \texttt{n} columns in \texttt{labels} refer to node labels; additional columns represent edge labels.
\end{itemize}

\subsubsection{Pseudocode: Graph Construction}

\begin{verbatim}
procedure initialize_graph(adj, labels, n, m)
    ! Allocate adj(n, m)
    ! Allocate labels(2, n + m)  ! n node labels, m edge labels
    set all entries of adj to -1
    set all labels to -1
end procedure

procedure add_node_label(labels, node_index, type_code, index)
    labels(1, node_index) = type_code
    labels(2, node_index) = index
end procedure

procedure add_edge(adj, labels, from_node, edge_slot, to_node_label_index)
    adj(from_node, edge_slot) = to_node_label_index
end procedure

procedure add_edge_label(labels, edge_label_index, type_code, index)
    labels(1, edge_label_index) = type_code
    labels(2, edge_label_index) = index
end procedure
\end{verbatim}

\subsubsection{Key Properties and Advantages}

\begin{itemize}
  \item Fully flat structure; ideal for streaming, serialization, and WASM.
  \item Separation of node/edge data from topology.
  \item Labels act as foreign keys into canonical data structures like \texttt{DataTable} or \texttt{HashMap}.
  \item Easy to extend for directed/undirected graphs or hypergraphs by modifying adjacency encoding.
  \item Compatible with graph algorithms via simple indexing and vector traversal.
\end{itemize}

\subsubsection{Test Case: Simple Directed Graph}

Example: A directed graph with 3 nodes (A, B, C) and 2 directed edges: A→B and B→C.

\begin{verbatim}
initialize_graph(adj, labels, 3, 2)

add_node_label(labels, 1, 1, 5)  ! Node A points to row 5 of DataTable
add_node_label(labels, 2, 1, 8)  ! Node B points to row 8
add_node_label(labels, 3, 1, 9)  ! Node C points to row 9

add_edge(adj, labels, 1, 1, 2)   ! A → B
add_edge(adj, labels, 2, 1, 3)   ! B → C
\end{verbatim}

This structure can now be queried, visualized, or extended without reallocation or symbolic interpretation.

\end{document}
